\documentclass[12pt,a4paper]{report}

% ----------------------------
% Packages
% ----------------------------
\usepackage[utf8]{inputenc}
\usepackage[T1]{fontenc}
\usepackage{lmodern}
\usepackage{geometry}
\geometry{margin=1in}

\usepackage{setspace}
\onehalfspacing

\usepackage{amsmath, amssymb, amsthm}
\usepackage{graphicx}
\usepackage{booktabs}
\usepackage{longtable}
\usepackage{hyperref}
\usepackage{tikz}
\usepackage{algorithm2e}
\usepackage{listings}
\usepackage{xcolor}
\usepackage{glossaries}
\usepackage{biblatex}
\addbibresource{references.bib}

\hypersetup{
    colorlinks=true,
    linkcolor=blue,
    urlcolor=blue,
    citecolor=blue
}

\makeglossaries

% ----------------------------
% Glossary Entries
% ----------------------------
\newglossaryentry{wal}{
    name=Write-Ahead Log,
    description={A logging mechanism ensuring durability and atomicity by recording changes before applying them}
}

\newglossaryentry{lsn}{
    name=LSN,
    description={Log Sequence Number, a monotonically increasing identifier for log records}
}

% ----------------------------
% Theorem Environments
% ----------------------------
\newtheorem{theorem}{Theorem}[chapter]
\newtheorem{definition}{Definition}[chapter]
\newtheorem{lemma}{Lemma}[chapter]

% ----------------------------
% Listings Configuration
% ----------------------------
\lstset{
    language=C,
    basicstyle=\ttfamily\small,
    keywordstyle=\color{blue},
    commentstyle=\color{green!60!black},
    stringstyle=\color{red},
    breaklines=true,
    frame=single
}

% ----------------------------
% Title Page
% ----------------------------
\title{\Huge Write-Ahead Logging (WAL) \\ \Large A Systems-Level Report}
\author{Technical Systems Analysis}
\date{\today}

\begin{document}

\maketitle
\tableofcontents
\newpage

% ----------------------------
% Abstract
% ----------------------------
\begin{abstract}
Write-Ahead Logging (WAL) is a core recovery mechanism in database systems ensuring atomicity and durability under failure. This report presents theoretical foundations, implementation mechanics, mathematical modeling, recovery algorithms, and practical trade-offs in WAL-based systems.
\end{abstract}

% ----------------------------
\chapter{Introduction}

\section{Motivation}

In durable storage systems, failures are inevitable. WAL ensures that:

\begin{enumerate}
    \item No committed transaction is lost.
    \item No partial transaction corrupts the database.
\end{enumerate}

\section{Formal Definition}

\begin{definition}
A \gls{wal} protocol requires that before any data page is written to persistent storage, the corresponding log record must be flushed to stable storage.
\end{definition}

% ----------------------------
\chapter{Mathematical Model of WAL}

\section{State Representation}

Let:

\[
D = \{p_1, p_2, ..., p_n\}
\]

be the set of data pages and

\[
L = \{r_1, r_2, ..., r_m\}
\]

be the ordered log records.

Each record has a \gls{lsn}:

\[
LSN(r_i) = i
\]

\section{WAL Invariant}

\[
\forall p \in D,\quad \text{if } pageLSN(p) = k,\text{ then log records } r_1...r_k \text{ are durable}
\]

\begin{theorem}
If the WAL invariant holds, then recovery guarantees atomicity and durability.
\end{theorem}

\begin{proof}
Since all modifications are logged before persistence, recovery can replay committed operations and undo incomplete ones using log records. Thus no committed state is lost and no partial state survives.
\end{proof}

% ----------------------------
\chapter{WAL Architecture}

\section{High-Level Components}

\begin{itemize}
    \item Log Buffer
    \item Log File (Sequential)
    \item Buffer Pool
    \item Checkpoint Manager
\end{itemize}

\section{Architecture Diagram}

\begin{center}
\begin{tikzpicture}[node distance=2cm]
\node (app) [rectangle, draw] {Transaction};
\node (logbuf) [rectangle, draw, below of=app] {Log Buffer};
\node (disklog) [rectangle, draw, below of=logbuf] {Log File (Disk)};
\node (buffer) [rectangle, draw, right of=logbuf, xshift=4cm] {Buffer Pool};
\node (diskdata) [rectangle, draw, below of=buffer] {Data Files};

\draw[->] (app) -- (logbuf);
\draw[->] (logbuf) -- (disklog);
\draw[->] (app) -- (buffer);
\draw[->] (buffer) -- (diskdata);
\end{tikzpicture}
\end{center}

% ----------------------------
\chapter{ARIES Recovery Algorithm}

\section{Phases}

\begin{enumerate}
    \item Analysis
    \item Redo
    \item Undo
\end{enumerate}

\section{Pseudo-code}

\begin{algorithm}[H]
\SetAlgoLined
\KwResult{Recovered database}
Read last checkpoint\;
Perform Analysis phase\;
Redo all operations from smallest dirty LSN\;
Undo incomplete transactions in reverse order\;
\caption{ARIES Recovery}
\end{algorithm}

% ----------------------------
\chapter{Implementation Example}

\section{Log Record Structure}

\begin{lstlisting}
typedef struct {
    uint64_t lsn;
    uint32_t txn_id;
    uint32_t page_id;
    uint32_t prev_lsn;
    char data[256];
} LogRecord;
\end{lstlisting}

\section{Commit Protocol}

\begin{enumerate}
    \item Append commit record to log buffer.
    \item Flush log buffer to disk.
    \item Acknowledge commit.
\end{enumerate}

% ----------------------------
\chapter{Checkpointing}

\section{Purpose}

Checkpointing reduces recovery time by limiting log replay.

\section{Checkpoint Record}

\[
Checkpoint = \{ActiveTxns, DirtyPages\}
\]

\section{Comparison Table}

\begin{longtable}{@{}lll@{}}
\toprule
Strategy & Pros & Cons \\
\midrule
Sharp Checkpoint & Fast recovery & High pause time \\
Fuzzy Checkpoint & No pause & More complex recovery \\
\bottomrule
\end{longtable}

% ----------------------------
\chapter{Performance Considerations}

\section{Sequential I/O Advantage}

Log writes are sequential:

\[
Throughput \propto \frac{1}{Seek\ Time}
\]

Sequential writes minimize seek overhead.

\section{Group Commit}

Multiple transactions share a single flush:

\[
Latency_{avg} = \frac{FlushTime}{N}
\]

% ----------------------------
\chapter{Advanced Topics}

\section{WAL in Distributed Systems}

Used in:
\begin{itemize}
    \item Replication logs
    \item Consensus protocols
    \item Event sourcing
\end{itemize}

\section{Logical vs Physical Logging}

\begin{itemize}
    \item Physical: byte-level changes
    \item Logical: semantic operations
\end{itemize}

% ----------------------------
\chapter{Conclusion}

Write-Ahead Logging remains the foundational durability mechanism in modern storage engines. Its strength lies in sequential I/O, formal invariants, and deterministic recovery semantics.

% ----------------------------
\appendix

\chapter{Glossary}
\printglossaries

\chapter{Future Research Directions}

\begin{itemize}
    \item Persistent memory optimization
    \item Hybrid logical-physical logging
    \item WAL compression techniques
\end{itemize}

% ----------------------------
\printbibliography

\end{document}
